% Hand-drawn (not data-generated): the QDCVR pipeline and where the two
% mechanisms under evaluation sit inside it. Natural width ~ textwidth.
\begin{figure*}[t]
\centering
\resizebox{0.80\textwidth}{!}{%
\begin{tikzpicture}[
  font=\small,
  stage/.style={draw, rounded corners=2pt, minimum height=11mm, minimum width=21mm,
                align=center, fill=black!4, inner sep=3pt},
  core/.style={draw, rounded corners=2pt, minimum height=11mm, minimum width=21mm,
               align=center, fill=black!16, line width=1.0pt, inner sep=3pt},
  store/.style={draw, dashed, rounded corners=2pt, minimum height=8.5mm,
                align=center, fill=white, inner sep=4pt, font=\footnotesize},
  ar/.style={-{Stealth[length=2mm]}, thick},
]

% ── main pipeline ────────────────────────────────────────────────
\node[stage] (q)   at (0,0)     {Query};
\node[stage] (s0)  at (2.35,0)  {S0\\ intent +\\ rewrite};
\node[core]  (s1)  at (4.9,0)   {\textbf{S1}\\ domain\\ scoping};
\node[stage] (s2)  at (7.45,0)  {S2\\ two-stage\\ recall};
\node[stage] (s3)  at (10.0,0)  {S3\\ dedup\\ + $\tau$};
\node[core]  (s4)  at (12.55,0) {\textbf{S4}\\ content\\ adjudication};
\node[stage] (s5)  at (15.1,0)  {S5\\ tiering};
\node[stage] (s6)  at (17.65,0) {S6\\ answer +\\ blind spot};

\foreach \a/\b in {q/s0, s0/s1, s1/s2, s2/s3, s3/s4, s4/s5, s5/s6}
  \draw[ar] (\a) -- (\b);

% ── stores ───────────────────────────────────────────────────────
\node[store] (kb)  at (4.9,-1.9)  {knowledge bases (hierarchical)};
\node[store] (idx) at (10.0,-1.9) {BM25 + ChromaDB index (full-corpus chunks)};
\node[store] (exp) at (15.1,-1.9) {experience store ($P_0/P_1/P_2$)};

\draw[ar, dashed] (s1) -- (kb);
\draw[ar, dashed] (s2) -- (idx);
\draw[ar, dashed] (s4.south) to[out=-100,in=140] (idx.north east);
\draw[ar, dashed] (exp) -- (s5);

% ── read path annotation ─────────────────────────────────────────
\draw[ar, dotted, thick] (s4.south west) to[out=-115,in=35]
   node[above, font=\footnotesize, pos=0.35]
   {\texttt{kb\_doc\_read}} (idx.west);

% ── highlight the two mechanisms ─────────────────────────────────
\node[draw, rounded corners=3pt, fit=(s1), inner sep=2.2mm, line width=1.2pt] {};
\node[draw, rounded corners=3pt, fit=(s4), inner sep=2.2mm, line width=1.2pt] {};

% ── agent surface ────────────────────────────────────────────────
\node[store, minimum width=118mm] (mcp) at (8.8,-3.6)
  {every stage reachable through the agent-native surface: 94 MCP tools + 14 agent skills (adjudication reads \emph{text}, not embeddings)};
\draw[ar, dashed] (s6.south) to[out=-90,in=20] (mcp.north east);

\end{tikzpicture}}
\caption{The \sys{} pipeline. Two stages carry the method's claims:
\textbf{S1 domain scoping} bounds the search to a small set of knowledge bases
before any document is compared, and \textbf{S4 content adjudication} reads
candidate text and scores it $0$--$8$, with authority to override the recall
order (Eq.~\ref{eq:override}). Every stage is reachable through the same MCP
tool surface the agents use.}
\label{fig:pipeline}
\end{figure*}
